\documentclass[11pt]{article}
\usepackage[a4paper,margin=1in]{geometry}
\usepackage{amsmath,amssymb,amsthm}
\usepackage{hyperref}

\newcommand{\dd}{\mathrm{d}}
\newcommand{\cD}{\mathcal{D}}
\newcommand{\cH}{\mathcal{H}}
\newcommand{\cQ}{\mathcal{Q}}
\newcommand{\cJ}{\mathcal{J}}
\newcommand{\fa}{\mathfrak{a}}
\newcommand{\avg}[1]{\langle #1\rangle}
\newcommand{\nstar}[1]{\|#1\|_*}

\newtheorem{theorem}{Theorem}
\newtheorem{lemma}[theorem]{Lemma}
\newtheorem{corollary}[theorem]{Corollary}
\theoremstyle{remark}
\newtheorem{remark}[theorem]{Remark}

\title{The collisionless limit of the heat-flux closure\\
for the linearized Landau--Coulomb operator}
\author{Working note}
\date{17 November 2026}

\begin{document}
\sloppy
\maketitle

\noindent\emph{Summary.} For the like-particle Landau operator with Coulomb
interaction, in three velocity dimensions, the quasi-static closure
$\hat q=-i\,\mathrm{sgn}(k)\,\theta(k)\,\hat T$ satisfies
\[
  \theta(k)\;\longrightarrow\;3\sqrt{2/\pi}\qquad(|k|\to\infty),
\]
the same value as for the Dougherty operator, and $\tfrac32$ times the
one-dimensional Hammett--Perkins value. Along the way: $k\,n_1\to\sqrt{\pi/2}$,
the perpendicular temperature response satisfies $k\,T_{\perp1}\to0$, and the
closure is positive for all sufficiently large $|k|$.

The proof is by matching upper and lower variational bounds, with trial
functions built from the Airy-type boundary layer of width
$(d(v_\perp)/k)^{1/3}$, where $d$ is the local velocity diffusion
coefficient. The limit does not depend on $d$.

\emph{Level of rigour.} The argument is complete in structure. The
asymptotic evaluation of the trial functionals is by dominated convergence;
the dominating bounds are stated, and the lower-order terms are estimated by
scaling with the reason given, but not every such term is written out. The
proof uses the cited estimates for the Landau operator, one more than
before (Section~\ref{sec:prop}). No second reader has checked it.

\section{Setting and properties used}
\label{sec:prop}

Notation is that of \texttt{landau\_solvability.tex} and
\texttt{landau\_quasistatic\_limit.tex}. $M(v)=M_1(v_x)M_\perp(v_\perp)$ is
the unit Maxwellian, $M_1(0)=(2\pi)^{-1/2}$, $\avg{\cdot}_\perp$ is the
average in $v_\perp$, and $m=(1+|v|^2)^{1/2}$. Take $k>0$. $\cD$ is the
Dirichlet form of the linearized Landau--Coulomb operator and
$g=\sqrt M\,h$. Besides (C1)--(C3):

\begin{enumerate}
  \item[(C6)] \emph{Local form plus a small remainder.} Let
    \[
      \fa[h]=\int\Bigl[S^{ij}\partial_ig\,\partial_jg
      +S^{ij}\tfrac{v_i}{2}\tfrac{v_j}{2}\,g^2\Bigr]\dd v ,
    \]
    with $S(v)$ the diffusion matrix (the matrix $\sigma^{ij}$ of the cited
    papers, up to the normalization of $\cD$). For every $\eta>0$ there are
    $C_\eta$, $R_\eta$ with
    \[
      (1-\eta)\,\fa[h]-C_\eta\!\int_{|v|\le R_\eta}\!Mh^2
      \ \le\ \cD[h]\ \le\
      (1+\eta)\,\fa[h]+C_\eta\!\int_{|v|\le R_\eta}\!Mh^2 .
    \]
    This is $L=-A-K$ with $-\langle Ag,g\rangle=|g|_\sigma^2
    -\langle\partial_i\sigma^ig,g\rangle$ and the bound on the
    $\partial_i\sigma^i$ and $K$ terms in \cite[Lemma~7]{StrainGuo}
    (equivalently \cite[Lemma~2.6]{KGH}).
  \item[(C7)] \emph{The diffusion matrix.} $S$ is smooth, with eigenvalue
    $\lambda_1(|v|)$ along $v$ and $\lambda_2(|v|)$ on $v^\perp$,
    $\lambda_1\sim m^{-3}$, $\lambda_2\sim m^{-1}$, and
    $|\nabla S|\le C m^{-2}$ \cite[Lemma~4]{StrainGuo}. Hence
    \[
      S^{xx}(v)\to d(v_\perp):=\lambda_2(|v_\perp|),\qquad
      (S^{-1})_{xx}(v)\to\frac{1}{d(v_\perp)}\qquad(v_x\to0),
    \]
    with $c\,m^{-1}\le d\le C$, $|\nabla d|\le Cd$, $S^{xx}\le C$,
    $(S^{-1})_{xx}\le C\,m^3$, and $S^{x\perp}=O(|v_x|)$.
\end{enumerate}

\paragraph{The dual seminorm.} For a function $\chi$ put
\[
  \nstar{\chi}=\sup\bigl\{\avg{\chi f}/\cD[f]^{1/2}:\ f\in\cH,\ \Pi f=0\bigr\}.
\]
It is a seminorm; functions in $H$ have $\nstar{\cdot}=0$; if $\chi\perp H$
then $\nstar{\chi}^2=\avg{\chi,\Lambda^+\chi}$; and by (C2), (C3),
$\nstar{\chi}\le C\avg{m\,\chi^2}^{1/2}$.

\paragraph{Variational quantities.} For even $\varphi\in\cH_\infty$ let
$\cQ_k[\varphi]=\cD[\varphi]+k^2\nstar{v_x\varphi}^2$. For an even polynomial
$\sigma=\sigma(v_\perp)$ define
\[
  N_k(\sigma)=\sup\Bigl\{\frac{\avg{\sigma\varphi}^2}{\cQ_k[\varphi]}:\
  \varphi\in\cH_\infty\text{ even},\ \avg{v_x^2\varphi}=0\Bigr\}.
\]
By the earlier notes $N_k(1)=n_1$, attained at $g_e$, and
$\cQ_k[g_e-\varphi]=n_1-(2\avg{\varphi}-\cQ_k[\varphi])$ for admissible
$\varphi$.

\begin{lemma}[Two-sided bounds]\label{lem:two}
Let $\sigma=\sigma(v_\perp)$ be an even polynomial.
\begin{enumerate}
  \item[(a)] For every even $\varphi\in\cH_\infty$, with
    $\tau=\avg{v_x^2\varphi}$,
    \[
      N_k(\sigma)\ \ge\ \frac{(\avg{\sigma\varphi}-\tau\avg{\sigma})^2}
      {\cD[\varphi]+k^2\nstar{v_x\varphi}^2}.
    \]
  \item[(b)] For every odd $\psi\in\cH_\infty$ such that
    $\sigma-kv_x\psi$ is orthogonal to $1$ and $|v|^2$,
    \[
      N_k(\sigma)\ \le\ \cD[\psi]+\nstar{\sigma-kv_x\psi}^2 .
    \]
  \item[(c)] For admissible $e$, $\avg{\sigma e}^2\le N_k(\sigma)\,\cQ_k[e]$.
\end{enumerate}
\end{lemma}

\begin{proof}
(a) $\varphi-\tau$ is admissible, $\cD[\varphi-\tau]=\cD[\varphi]$, and
$\nstar{v_x(\varphi-\tau)}=\nstar{v_x\varphi}$ because $v_x\in H$.
(b) For admissible $\varphi$ and $\chi=\sigma-kv_x\psi$,
$\avg{\sigma\varphi}=\avg{\chi\varphi}+k\avg{v_x\varphi\,\psi}$. Since
$\chi\perp\Pi\varphi$ and $v_x\varphi\perp\Pi\psi$, the two terms are at most
$\nstar{\chi}\cD[\varphi]^{1/2}$ and $k\nstar{v_x\varphi}\cD[\psi]^{1/2}$;
apply the Cauchy--Schwarz inequality in $\mathbb{R}^2$. (c) is the
definition.
\end{proof}

\section{The boundary-layer profile}

\begin{lemma}\label{lem:airy}
Let $\gamma(u)=\int_0^\infty e^{-iut-t^3/3}\,\dd t=F(u)+iW(u)$. Then:
\begin{enumerate}
  \item[(a)] $iu\gamma-\gamma''=1$; that is, $-uW-F''=1$ and $uF-W''=0$.
  \item[(b)] $F$ is even and $W$ odd; as $|u|\to\infty$,
    $F=-2u^{-4}+O(u^{-7})$, $W=-u^{-1}+O(u^{-4})$, with corresponding
    bounds for derivatives.
  \item[(c)] $\int F\,\dd u=\pi$, and
    $\int_u^\infty tF(t)\,\dd t=-W'(u)$.
  \item[(d)] $\int F'^2\,\dd u+\int W'^2\,\dd u=\pi$.
\end{enumerate}
\end{lemma}

\begin{proof}
(a) Differentiate under the integral and integrate by parts:
$\gamma''-iu\gamma=\int_0^\infty\frac{\dd}{\dd t}\bigl(e^{-iut}\bigr)
e^{-t^3/3}\dd t-\int_0^\infty t^2e^{-iut-t^3/3}\dd t=-1$. (b) Parity is
immediate; the expansion is repeated integration by parts, with
$a(t)=e^{-t^3/3}$, $a(0)=1$, $a'(0)=a''(0)=0$, $a'''(0)=-2$:
$\gamma=(iu)^{-1}-2(iu)^{-4}+\dots$. (c) $\int\gamma\,\dd u=\pi$ by Fourier
inversion on the half-line. From $W''=uF$ and $W'\to0$ at both ends,
$W'(u)=-\int_u^\infty tF$. (d) Multiply the first equation in (a) by $F$ and
the second by $W$, integrate, and add: $\int F'^2+\int W'^2=\int F=\pi$; the
boundary terms vanish by (b).
\end{proof}

These identities were also checked numerically.

\section{A sharp upper bound for the dual seminorm}

\begin{lemma}\label{lem:flux}
Let $\chi_k$ be smooth functions of the form
$\chi_k=-M^{-1}\partial_{v_x}(MJ_k)$, with $J_k$ decaying at infinity, and
put $\cJ_k=\int MJ_k^2\,(S^{-1})_{xx}\,\dd v$. Suppose that, as $k\to\infty$,
\[
  \int M|J_k|\,\dd v=o(\cJ_k^{1/2}),\qquad
  \int MJ_k^2\,v_x^2\,m\,\dd v=o(\cJ_k).
\]
Then $\nstar{\chi_k}^2\le(1+o(1))\,\cJ_k$.
\end{lemma}

\begin{proof}
Write $s_k=\nstar{\chi_k}^2$ and assume $s_k>\cJ_k$ along a subsequence,
otherwise there is nothing to prove. By the Riesz theorem there is
$f_k\in\cH$, $\Pi f_k=0$, with $\cD(f_k,\zeta)=\avg{\chi_k\zeta}$ for
$\zeta\perp H$ and $\cD[f_k]=\avg{\chi_kf_k}=s_k$. With $g_k=\sqrt M f_k$,
\[
  s_k=\int MJ_k\,\partial_xf_k
     =\int\sqrt M\,J_k\,\partial_xg_k+\int MJ_k\,\tfrac{v_x}{2}\,f_k .
\]
The first term is at most
$\cJ_k^{1/2}\bigl(\int S^{ij}\partial_ig_k\partial_jg_k\bigr)^{1/2}
\le\cJ_k^{1/2}\fa[f_k]^{1/2}$. The second is at most
$\bigl(\int MJ_k^2v_x^2m\bigr)^{1/2}\avg{m^{-1}f_k^2}^{1/2}
=o(\cJ_k^{1/2})\,s_k^{1/2}$ by (C2), (C3).

By (C6), $\fa[f_k]\le(1-\eta)^{-1}\bigl(s_k+C_\eta\int_{|v|\le R_\eta}Mf_k^2\bigr)$.
Let $\tilde f_k=f_k/s_k^{1/2}$, bounded in $\cH^\perp$. For smooth compactly
supported $\zeta$,
$|\cD(\tilde f_k,\zeta-\Pi\zeta)|=|\int MJ_k\partial_x(\zeta-\Pi\zeta)|/s_k^{1/2}\to0$
by the two hypotheses on $J_k$ (the second handles the term $2v_x$ coming
from $\partial_x\Pi\zeta$), so $\tilde f_k\to0$ weakly in
$\cH^\perp$. The norm of $\cH$ controls the $H^1$ norm on bounded sets, by
(C2), so by Rellich's theorem $\int_{|v|\le R_\eta}M\tilde f_k^2\to0$.
Therefore $\fa[f_k]\le(1-\eta)^{-1}(1+o(1))\,s_k$, and
\[
  s_k\le\cJ_k^{1/2}s_k^{1/2}\bigl((1-\eta)^{-1/2}+o(1)\bigr).
\]
Since $\eta$ is arbitrary, $s_k\le(1+o(1))\cJ_k$.
\end{proof}

\section{The theorem}

\begin{theorem}\label{thm:main}
For every even polynomial $\sigma=\sigma(v_\perp)$,
\[
  \lim_{k\to\infty}k\,N_k(\sigma)=\pi\,M_1(0)\,\avg{\sigma^2}_\perp .
\]
\end{theorem}

Throughout, $\delta=\delta(v_\perp)=(d(v_\perp)/k)^{1/3}$ and
$u=v_x/\delta$; by (C7), $c\,k^{-1/3}m^{-1/3}\le\delta\le Ck^{-1/3}$ and
$|\nabla_\perp\delta|\le C\delta$. Integrals over $v_x$ are rescaled to $u$,
and limits are taken under the integral in $(u,v_\perp)$ by dominated
convergence; the Gaussian $M_\perp$ and the decay in Lemma~\ref{lem:airy}(b)
supply the domination.

\begin{proof}[Proof of the lower bound]
Let $\varphi_k(v)=\sigma(v_\perp)\,\delta^{-1}F(v_x/\delta)$ and apply
Lemma~\ref{lem:two}(a).

\emph{Numerator.}
$\avg{\sigma\varphi_k}=\int M_\perp\sigma^2\int M_1(\delta u)F(u)\,\dd u\,\dd v_\perp
\to\pi M_1(0)\avg{\sigma^2}_\perp$, and
$\tau_k=\int M_\perp\sigma\,\delta^2\int u^2F(u)M_1(\delta u)\,\dd u\,\dd v_\perp
=O(k^{-2/3})$, since $u^2F$ is integrable.

\emph{$\cD[\varphi_k]$.} The leading part of $\fa[\varphi_k]$ is
\begin{align*}
  \int MS^{xx}|\partial_x\varphi_k|^2
  &=\int M_\perp\sigma^2\delta^{-3}\int S^{xx}(\delta u,v_\perp)M_1(\delta u)
   F'(u)^2\,\dd u\,\dd v_\perp \\
  &=k\,M_1(0)\avg{\sigma^2}_\perp\!\int F'^2\,(1+o(1)),
\end{align*}
because $\delta^{-3}d=k$. The remaining terms of $\fa$ are of lower order:
the perpendicular derivatives of $\varphi_k$ are $O(\delta^{-1})$ on a layer
of width $\delta$ and contribute $O(k^{1/3})$; the cross terms contribute
$O(k^{2/3})$ by the Cauchy--Schwarz inequality; the terms without
derivatives, and those in which the derivative falls on $\sqrt M$, are
$O(k^{1/3})$ and $O(k^{2/3})$. Likewise $\int_{|v|\le R}M\varphi_k^2=O(k^{1/3})$.
By (C6), $\cD[\varphi_k]\le(1+\eta)\,kM_1(0)\avg{\sigma^2}_\perp\int F'^2\,(1+o(1))$.

\emph{$k^2\nstar{v_x\varphi_k}^2$.} $\chi_k=v_x\varphi_k=\sigma\,uF(u)$ is odd
in $v_x$, so $\chi_k=-M^{-1}\partial_x(MJ_k)$ with
\[
  J_k(v)=\frac{1}{M_1(v_x)}\int_{v_x}^\infty\chi_k\,M_1\,\dd s
        =\sigma\,\delta\,j_\delta(u),\qquad
  j_\delta(u)=\frac{1}{M_1(\delta u)}\int_u^\infty tF(t)M_1(\delta t)\,\dd t .
\]
Oddness gives the same formula with $-\int_{-\infty}^{u}$, so
$|j_\delta(u)|\le C(1+u^2)^{-1}$ uniformly, and $j_\delta\to-W'$ by
Lemma~\ref{lem:airy}(c). Hence, using $(S^{-1})_{xx}\to1/d$ and
$\delta^3/d=1/k$,
\[
  \cJ_k=\int M_\perp\sigma^2\delta^3\!\int M_1(\delta u)(S^{-1})_{xx}\,
  j_\delta(u)^2\,\dd u\,\dd v_\perp
  =\frac1k\,M_1(0)\avg{\sigma^2}_\perp\!\int W'^2\,(1+o(1)).
\]
(The factor $m^3$ in the bound for $(S^{-1})_{xx}$ is harmless: for
$|v_x|\le1$ it is bounded by a polynomial in $v_\perp$, and for $|v_x|>1$ the
contribution is $O(\delta^6)$.) The hypotheses of Lemma~\ref{lem:flux} hold:
$\int M|J_k|=O(\delta^2)=o(k^{-1/2})$, and
$\int MJ_k^2v_x^2m=O(\delta^5)=o(k^{-1})$. So
$k^2\nstar{v_x\varphi_k}^2\le kM_1(0)\avg{\sigma^2}_\perp\int W'^2\,(1+o(1))$.

\emph{Conclusion.} By Lemma~\ref{lem:airy}(d),
$\cQ_k[\varphi_k]\le(1+\eta)\,\pi M_1(0)\avg{\sigma^2}_\perp k\,(1+o(1))$, so
$\liminf kN_k(\sigma)\ge\pi M_1(0)\avg{\sigma^2}_\perp/(1+\eta)$ for every
$\eta>0$.
\end{proof}

\begin{proof}[Proof of the upper bound]
Let
$\psi_k=-\sigma\,(k\delta)^{-1}W(v_x/\delta)+k^{-1}(\alpha_kv_x+\beta_kv_x|v|^2)$
and $\chi_k=\sigma-kv_x\psi_k$. By Lemma~\ref{lem:airy}(a),
$\sigma(1+uW)=-\sigma F''(u)$, so
\[
  \chi_k=-\sigma\,F''(v_x/\delta)-\alpha_kv_x^2-\beta_kv_x^2|v|^2 .
\]
Two integrations by parts in $v_x$ give
$\int F''(v_x/\delta)M_1\,\dd v_x=\delta^2\!\int F(v_x/\delta)M_1''\,\dd v_x
=O(\delta^3)$, so the moments of $-\sigma F''$ against $1$ and $|v|^2$ are
$O(k^{-1})$, and $\alpha_k,\beta_k=O(k^{-1})$ can be chosen to make $\chi_k$
orthogonal to $1$ and $|v|^2$. Apply Lemma~\ref{lem:two}(b).

\emph{$\cD[\psi_k]$.} The leading part of $\fa[\psi_k]$ is
\begin{align*}
  \int MS^{xx}|\partial_x\psi_k|^2
  &=\frac{1}{k^2}\int M_\perp\sigma^2\delta^{-3}\!\int S^{xx}M_1(\delta u)
   W'(u)^2\,\dd u\,\dd v_\perp \\
  &=\frac1k\,M_1(0)\avg{\sigma^2}_\perp\!\int W'^2\,(1+o(1)).
\end{align*}
$W'^2$ is integrable. The slowly decaying tail $W\sim-1/u$ means
$\psi_k\approx\sigma/(kv_x)$ away from the layer; there every term of
$\fa[\psi_k]$ is $O(k^{-2})$. In the layer the remaining terms are
$O(k^{-2}\delta^{-1})=O(k^{-5/3})$, as is $\int_{|v|\le R}M\psi_k^2$, since
$W^2$ is integrable. The corrections with $\alpha_k,\beta_k$ contribute
$O(k^{-4})$. So $\cD[\psi_k]\le(1+\eta)\,k^{-1}M_1(0)\avg{\sigma^2}_\perp
\int W'^2\,(1+o(1))$.

\emph{$\nstar{\chi_k}$.} Let
$\mu_k(v_\perp)=-\sigma\int F''(s/\delta)M_1(s)\,\dd s=O(\delta^3)$ and
split $\chi_k=\chi_k^0+r_k$, where $\chi_k^0=-\sigma F''(v_x/\delta)-\mu_k$
has zero $M_1$-mean in $v_x$ for every $v_\perp$, and
$r_k=\mu_k-\alpha_kv_x^2-\beta_kv_x^2|v|^2$. Then
$\nstar{r_k}\le C\avg{m\,r_k^2}^{1/2}=O(k^{-1})$. For $\chi_k^0$,
\[
  J_k^0(v)=\frac{1}{M_1(v_x)}\int_{v_x}^\infty\chi_k^0M_1\,\dd s
  =\sigma\,\delta\,F'(u)\,(1+o(1))+O(\delta^3),
\]
bounded by $C\sigma\delta(1+|u|)^{-5}+C\delta^3$ using the zero mean on
either side. Hence
$\cJ_k^0=k^{-1}M_1(0)\avg{\sigma^2}_\perp\int F'^2\,(1+o(1))$, the hypotheses
of Lemma~\ref{lem:flux} hold as before, and
\[
  \nstar{\chi_k}^2\le\bigl(\nstar{\chi_k^0}+\nstar{r_k}\bigr)^2
  \le\frac1k\,M_1(0)\avg{\sigma^2}_\perp\!\int F'^2\,(1+o(1)).
\]

\emph{Conclusion.} By Lemma~\ref{lem:airy}(d),
$\limsup kN_k(\sigma)\le(1+\eta)\,\pi M_1(0)\avg{\sigma^2}_\perp$ for every
$\eta>0$.
\end{proof}

\section{Consequences}

\begin{corollary}\label{cor:theta}
For the linearized Landau--Coulomb operator, as $k\to\infty$,
\[
  k\,n_1\to\sqrt{\pi/2},\qquad
  k\,T_{\perp1}\to0,\qquad
  -k\,T_1\to\tfrac13\sqrt{\pi/2},\qquad
  \theta(k)\to3\sqrt{2/\pi}.
\]
In particular $\chi(k)>0$ for all sufficiently large $k$, and
$\chi(k)\sim3\sqrt{2/\pi}\,/k$.
\end{corollary}

\begin{proof}
$n_1=N_k(1)$ and $\pi M_1(0)=\sqrt{\pi/2}$. Let $\varphi_k$ be the trial
function of the lower bound with $\sigma=1$,
$\varphi_k'=\varphi_k-\tau_k$, and $t_k=\avg{\varphi_k'}/\cQ_k[\varphi_k']$.
The two bounds of Theorem~\ref{thm:main} force
$\cQ_k[\varphi_k']/k\to\pi M_1(0)$, so $kt_k\to1$, and by the gap identity,
\[
  \cQ_k[g_e-t_k\varphi_k']
  =n_1-\frac{\avg{\varphi_k'}^2}{\cQ_k[\varphi_k']}=o(k^{-1}).
\]
By Lemma~\ref{lem:two}(c) and Theorem~\ref{thm:main} for
$\sigma=v_\perp^2$,
$|\avg{v_\perp^2(g_e-t_k\varphi_k')}|\le
\bigl(N_k(v_\perp^2)\,\cQ_k[g_e-t_k\varphi_k']\bigr)^{1/2}=o(k^{-1})$.
Also $\avg{v_\perp^2\varphi_k'}\to\pi M_1(0)\avg{v_\perp^2}_\perp
=2\sqrt{\pi/2}$. Hence $kX=k\avg{v_\perp^2g_e}\to2\sqrt{\pi/2}$. Since
$X=2(n_1+T_{\perp1})$ and $-T_1=n_1-\tfrac13X$, the remaining limits follow,
and $\theta=-1/(kT_1)$.
\end{proof}

\begin{remark}[Universality]
The proof uses only (C2), (C3), (C6), (C7): coercivity, a local diffusion
form plus a relatively compact remainder, and a continuous positive
diffusion coefficient at $v_x=0$. The limit depends on none of the
constants. For the Dougherty operator $S=\nu I$ and the theorem recovers the
value obtained from the closed form.
\end{remark}

\begin{remark}[What this does and does not give]
The convergence has no rate, so there is no explicit $k_1$ beyond which
positivity holds. Together with the earlier result, positivity is proved for
$k\lambda<0.89$ and for $k\lambda>k_1$ with $k_1$ unknown; the intermediate
range is covered only numerically. A rate would follow from quantitative
versions of the dominated-convergence steps and of the compactness step in
Lemma~\ref{lem:flux}, which is the one non-constructive point.
\end{remark}

\section{Checks}

\begin{itemize}
  \item Lemma~\ref{lem:airy} was verified numerically: the residual of the
    differential equation is below $10^{-4}$, $\int F=3.14160$,
    $\int F'^2+\int W'^2=3.1415927$, $u^4F\to-2.0000$, $uW\to-1.0000$.
  \item The limit agrees with the Dougherty closed form, where
    $\theta_3\to3\sqrt{2/\pi}=2.3937$ and the computed values are $2.3825$
    and $2.3899$ at $\nu/k=0.003$ and $0.001$.
  \item The Landau limit itself is not checked numerically: the polynomial
    Galerkin method does not reach large $k$. At $k\lambda=5$ it gives
    $-T_1/n_1=0.396$, against the limiting value $\tfrac13$, and
    $T_{\perp1}/n_1=-0.094$, against $0$.
\end{itemize}

\begin{thebibliography}{9}
\bibitem{StrainGuo} R.~M. Strain, Y.~Guo, \emph{Exponential decay for soft
potentials near Maxwellian}, Arch.\ Ration.\ Mech.\ Anal.\ \textbf{187}
(2008) 287--339 (Lemmas 4 and 7 read).
\bibitem{KGH} J.~Kim, Y.~Guo, H.~J. Hwang, \emph{A $L^2$ to $L^\infty$
approach for the Landau equation}, arXiv:1610.05346 (Lemmas 2.3, 2.5, 2.6
read; these restate lemmas of Guo, Comm.\ Math.\ Phys.\ 231 (2002), which was
not consulted directly).
\end{thebibliography}

\end{document}
